\section{Matter-response theorem driven by the physical $n=2$ mode}
\label{sec:bhsm-matter-response}

The metric-response theorem fixes the scalar metric perturbation generated by
the reduced BHSM phase-space state $X_2=(q_2,\Pi_2)^T$.  Matter and radiation
then respond through covariant stress-energy conservation on that induced
metric.  No additional topographic degree of freedom is introduced at this
stage.

\subsection{Forced conservation system}

Let
\begin{equation}
Y_m
=
\begin{pmatrix}
\delta_m\\
v_m\\
\delta_r\\
v_r
\end{pmatrix}
\end{equation}
denote the linear dust/radiation response variables in a fixed scalar gauge.
For a specified background and scalar metric response $h_2$, covariant
conservation has the linear form
\begin{equation}
\frac{dY_m}{dN}
=
\mathbb A_m(N)Y_m
+
\mathbb S_m(N)h_2(N),
\label{eq:matter-forced-system}
\end{equation}
where $N=\ln a$.  Substituting the metric-response theorem,
\begin{equation}
h_2
=
h_{\rm src}
+
\mathcal R_{g2}X_2,
\end{equation}
gives
\begin{equation}
\frac{dY_m}{dN}
=
\mathbb A_mY_m
+
\mathbb S_m h_{\rm src}
+
\mathbb S_m\mathcal R_{g2}X_2 .
\label{eq:matter-forced-by-q2}
\end{equation}

\subsection{Matter-response theorem}

\paragraph{Theorem 3 (unique matter response).}
Assume that the linear matter conservation problem is well posed on the
chosen redshift interval, with fundamental matrix
$\mathbb U_m(N,N_i)$ for the homogeneous system
$Y_m'=\mathbb A_mY_m$.  Then the response sourced by the physical $n=2$
BHSM mode is
\begin{equation}
\boxed{
Y_{m,2}(N)
=
\mathbb U_m(N,N_i)Y_{m,2}(N_i)
+
\int_{N_i}^{N}
\mathbb U_m(N,N')
\mathbb S_m(N')
\mathcal R_{g2}(N')X_2(N')\,dN' .
}
\label{eq:bhsm-matter-response}
\end{equation}
For a fixed physical initial-data prescription, this defines a linear matter
response operator
\begin{equation}
\boxed{
Y_{m,2}(N)=\mathcal R_{m2}(N)X_2
}
\label{eq:Rm2}
\end{equation}
in the enlarged sense that $\mathcal R_{m2}$ includes the causal history
integral generated by the reduced mode.

\paragraph{Proof.}
Equation~\eqref{eq:matter-forced-by-q2} is an inhomogeneous first-order linear
system.  Variation of constants yields
Eq.~\eqref{eq:bhsm-matter-response}.  Once the initial-data prescription is
fixed as part of the physical branch, every term is linear in the retained
mode data and its causal history.  Hence the response can be represented by
the linear history operator $\mathcal R_{m2}$. \hfill$\square$

\subsection{Adiabatic branch and absence of extra free data}

For the single-mode cosmological branch, the matter initial conditions are
not treated as arbitrary anisotropic fit functions.  The reference branch
uses one of the following explicitly declared choices:

\begin{enumerate}
\item an adiabatic matter/radiation seed fixed by the background branch;
\item a vanishing anisotropic matter seed, so that all directional response is
      induced by the BHSM topographic mode;
\item an independently measured environmental source, written separately from
      the one-mode prediction.
\end{enumerate}

Under the second choice, the directional matter response is completely
induced by $X_2$.  Under the first choice, the isotropic growing mode and the
directional $n=2$ response remain separately identifiable.

\subsection{Harmonic preservation}

Because $\mathbb A_m$, $\mathbb S_m$, and $\mathcal R_{g2}$ are rotationally
covariant on the closed background, a pure physical $n=2$ source cannot
generate a different scalar harmonic at linear order in the absence of an
additional anisotropic source.

\paragraph{Corollary 3.1.}
For the pure $A_{4i}$ component,
\begin{align}
\delta_m(z,\hat{\mathbf n})
&=
\delta_{m,2}(z)
(\hat{\mathbf n}\!\cdot\!\hat{\mathbf p}),
\\
v_m(z,\hat{\mathbf n})
&=
v_{m,2}(z)
(\hat{\mathbf n}\!\cdot\!\hat{\mathbf p}),
\\
\delta_r(z,\hat{\mathbf n})
&=
\delta_{r,2}(z)
(\hat{\mathbf n}\!\cdot\!\hat{\mathbf p}),
\\
v_r(z,\hat{\mathbf n})
&=
v_{r,2}(z)
(\hat{\mathbf n}\!\cdot\!\hat{\mathbf p}).
\end{align}
Thus the leading matter, metric, distance, growth, and lensing responses share
the same one-mode axis.

\subsection{Growth-response functional}

Let $D(N)$ denote the isotropic reference growth factor and let
$\delta D_2(N)$ denote the $n=2$ directional growth response obtained from
$\mathcal R_{m2}$.  Then, to linear order,
\begin{equation}
\frac{\delta(f\sigma_8)}{f\sigma_8}
=
\left[
\frac{\delta f_2}{f}
+
\frac{\delta D_2}{D}
\right]
(\hat{\mathbf n}\!\cdot\!\hat{\mathbf p}),
\label{eq:fs8-response}
\end{equation}
with
\begin{equation}
\delta f_2
=
\frac{d}{dN}
\left(
\frac{\delta D_2}{D}
\right).
\end{equation}
Equation~\eqref{eq:fs8-response} is therefore not an independent ansatz:
once $\mathcal R_{m2}$ is evaluated, the directional growth kernel is fixed.

\subsection{Combined physical transfer operator}

Combining the metric and matter responses gives the physical linear response
operator
\begin{equation}
\boxed{
\mathcal R_2
=
\begin{pmatrix}
\mathcal R_{g2}\\
\mathcal R_{m2}
\end{pmatrix}.
}
\label{eq:combined-R2}
\end{equation}
For any scalar linear observable,
\begin{equation}
\boxed{
\mathcal A_{\mathcal O}(z)
=
\mathcal L_{\mathcal O,z}
\mathcal R_2(z)X_2 .
}
\label{eq:final-observable-kernel}
\end{equation}
This is the closed operator-level bridge from the BHSM action-owned mode to
late-time observables.

\subsection{Numerical closure object}

After Theorems 1--3, the remaining numerical task is no longer the discovery
of a missing degree of freedom.  It is the evaluation of the already-defined
operators
\begin{equation}
\mathcal G_2,\quad
\mathcal M_2^2,\quad
\mathcal J_2,\quad
\mathcal R_{g2},\quad
\mathcal R_{m2}
\end{equation}
on the frozen reference branch.  Those quantities determine the observable
amplitudes in Eq.~\eqref{eq:final-observable-kernel}.

The reference-branch curves quoted later in the manuscript should therefore be
read as candidate numerical evaluations of this operator chain, pending exact
microscopic normalization of every reduced block.
